\documentclass[10pt,a4paper]{amsart}
\usepackage{amsmath,amssymb,amsthm,graphicx}
% Retained preamble for the Stacks Project chapter morphisms.tex (Morphisms of Schemes).
% The chapter file pinned-original-source.tex carries only \input{preamble} on its first line, so no
% preamble text lives inside the chapter itself.  The Stacks Project uses one single shared preamble for
% every chapter; the snapshot stages it verbatim in sources/stacks-chow/preamble.tex.  The block below is
% copied mechanically, line for line, from that file's lines 14-18 (the Xy-pic setup) and lines 154-217
% (the native theorem declarations together with the mathematical macros).  Nothing was added, removed or
% redefined.  The wrapper class and the amsmath/amssymb/amsthm/graphicx base packages are supplied by the
% assembler, so the class/\externaldocument part of the shared preamble is not repeated here.
% For commutative diagrams we use Xy-pic
\usepackage[all]{xy}

% We use 2cell for 2-commutative diagrams.
\xyoption{2cell}
% Theorem environments.
%
\theoremstyle{plain}
\newtheorem{theorem}[subsection]{Theorem}
\newtheorem{proposition}[subsection]{Proposition}
\newtheorem{lemma}[subsection]{Lemma}

\theoremstyle{definition}
\newtheorem{definition}[subsection]{Definition}
\newtheorem{example}[subsection]{Example}
\newtheorem{exercise}[subsection]{Exercise}
\newtheorem{situation}[subsection]{Situation}

\theoremstyle{remark}
\newtheorem{remark}[subsection]{Remark}
\newtheorem{remarks}[subsection]{Remarks}

\numberwithin{equation}{subsection}

% Macros
%
\def\lim{\mathop{\mathrm{lim}}\nolimits}
\def\colim{\mathop{\mathrm{colim}}\nolimits}
\def\Spec{\mathop{\mathrm{Spec}}}
\def\Hom{\mathop{\mathrm{Hom}}\nolimits}
\def\Ext{\mathop{\mathrm{Ext}}\nolimits}
\def\SheafHom{\mathop{\mathcal{H}\!\mathit{om}}\nolimits}
\def\SheafExt{\mathop{\mathcal{E}\!\mathit{xt}}\nolimits}
\def\Sch{\mathit{Sch}}
\def\Mor{\mathop{\mathrm{Mor}}\nolimits}
\def\Ob{\mathop{\mathrm{Ob}}\nolimits}
\def\Sh{\mathop{\mathit{Sh}}\nolimits}
\def\NL{\mathop{N\!L}\nolimits}
\def\CH{\mathop{\mathrm{CH}}\nolimits}
\def\proetale{{pro\text{-}\acute{e}tale}}
\def\etale{{\acute{e}tale}}
\def\QCoh{\mathit{QCoh}}
\def\Ker{\mathop{\mathrm{Ker}}}
\def\Im{\mathop{\mathrm{Im}}}
\def\Coker{\mathop{\mathrm{Coker}}}
\def\Coim{\mathop{\mathrm{Coim}}}

% Boxtimes
%
\DeclareMathSymbol{\boxtimes}{\mathbin}{AMSa}{"02}

%
% Macros for moduli stacks/spaces
%
\def\QCohstack{\mathcal{QC}\!\mathit{oh}}
\def\Cohstack{\mathcal{C}\!\mathit{oh}}
\def\Spacesstack{\mathcal{S}\!\mathit{paces}}
\def\Quotfunctor{\mathrm{Quot}}
\def\Hilbfunctor{\mathrm{Hilb}}
\def\Curvesstack{\mathcal{C}\!\mathit{urves}}
\def\Polarizedstack{\mathcal{P}\!\mathit{olarized}}
\def\Complexesstack{\mathcal{C}\!\mathit{omplexes}}
% \Pic is the operator that assigns to X its picard group, usage \Pic(X)
% \Picardstack_{X/B} denotes the Picard stack of X over B
% \Picardfunctor_{X/B} denotes the Picard functor of X over B
\def\Pic{\mathop{\mathrm{Pic}}\nolimits}
\def\Picardstack{\mathcal{P}\!\mathit{ic}}
\def\Picardfunctor{\mathrm{Pic}}
\def\Deformationcategory{\mathcal{D}\!\mathit{ef}}

% --- appended 2026-10-12 (Chief): the project own definitions of the reference/slogan/history
% display environments, copied mechanically from the shared Stacks preamble lines 8-12
% (\usepackage{verbatim} plus the three \newenvironment lines, which the chapter preamble keeps
% verbatim from the project preamble). Nothing added, removed or redefined: these environments
% carry no mathematical content and the project itself typesets them as comments. Without these
% three lines a counted statement whose own text uses them fails with Environment slogan undefined.
\usepackage{verbatim}
\newenvironment{reference}{\comment}{\endcomment}
\newenvironment{slogan}{\comment}{\endcomment}
\newenvironment{history}{\comment}{\endcomment}
\usepackage{cleveref}
\crefname{originallemma}{original lemma}{original lemmas}
\Crefname{originallemma}{Original lemma}{Original lemmas}
\makeatletter
\newlabel{morphisms-lemma-distinct-local-rings}{{02NF}{source}{}{sourcepointer-1}{}}
\newlabel{morphisms-lemma-distinct-local-rings@cref}{{[originallemma][0][]02NF}{[1][0][]source}{}{}{}}
\makeatother
\crefname{originallemma}{original lemma}{original lemmas}
\Crefname{originallemma}{Original lemma}{Original lemmas}
\makeatletter
\newlabel{schemes-equation-canonical-morphism}{{02NA}{source}{}{sourcepointer-2}{}}
\newlabel{schemes-equation-canonical-morphism@cref}{{[originallemma][0][]02NA}{[1][0][]source}{}{}{}}
\makeatother
\crefname{originallemma}{original lemma}{original lemmas}
\Crefname{originallemma}{Original lemma}{Original lemmas}
\makeatletter
\newlabel{schemes-lemma-where-are-they-equal}{{01KM}{source}{}{sourcepointer-3}{}}
\newlabel{schemes-lemma-where-are-they-equal@cref}{{[originallemma][0][]01KM}{[1][0][]source}{}{}{}}
\makeatother
\crefname{originallemma}{original lemma}{original lemmas}
\Crefname{originallemma}{Original lemma}{Original lemmas}
\makeatletter
\newlabel{schemes-lemma-specialize-points}{{01J7}{source}{}{sourcepointer-4}{}}
\newlabel{schemes-lemma-specialize-points@cref}{{[originallemma][0][]01J7}{[1][0][]source}{}{}{}}
\makeatother
\title{Distinct irreducible closed subsets of an integral separated scheme have distinct local rings in the function field}
\author{The Stacks Project Authors}
\date{stacks-project:morphisms.tex:local-snapshot-20261012; GNU FDL 1.2 or later}
\begin{document}
\maketitle
\setcounter{subsection}{1}
\noindent\textit{Editorial note.} All delivered mathematical text is the author's own, copied from the pinned Stacks Project chapter file morphisms.tex; no statement or proof step was rewritten, reordered or removed.\par\medskip
\noindent\textit{Editorial note.} The chapter file pinned-original-source.tex carries only \textbackslash{}input\{preamble\} on its first line, so no preamble text lives inside the chapter. The retained file sources/stacks-morphisms/retained-preamble.tex is a mechanical, line-for-line copy of the project's single shared preamble as staged in the snapshot. Nothing was added, removed or redefined, and the excerpt serves only as preamble, never as a span coordinate.\par\medskip
\noindent\textit{Editorial note.} Both counted spans start at the \textbackslash{}begin line of their environment, so the statement and the author's proof are delivered whole; surrounding prose of the section is deliberately outside the counted spans and is carried as standard prerequisites instead.\par\medskip
\noindent\textit{Editorial note.} The citation fields (source\_id, source\_version, authors, source\_url, license.name, source\_title) are identical to those of every other card this lane has produced from this chapter, so the batch selector treats them as one source.\par\medskip
\noindent\textit{Allowed standard prerequisites (editorial).}\begin{itemize}
\item Local rings of an integral scheme are local subrings of the function field. Lemma \textbackslash{}ref\{lemma-integral-scheme-rational-functions\}, cited at proof lines 12752-12754. Makes the two local rings comparable as subrings of one field.
\item A unique morphism to the scheme from the spectrum of the local ring, compatible with the canonical map from the function field. Schemes, Equation \textbackslash{}ref\{schemes-equation-canonical-morphism\} and Schemes, Lemma \textbackslash{}ref\{schemes-lemma-where-are-they-equal\}, cited at proof lines 12773-12785. Separatedness forces that morphism to be unique.
\item Specialisation of points read off from a morphism of spectra. Schemes, Lemma \textbackslash{}ref\{schemes-lemma-specialize-points\}, cited at proof lines 12799-12801. Turns the hypothetical inclusion into a containment of the closures.
\end{itemize}\medskip
\ifdefined\hypertarget\hypertarget{sourcepointer-1}{}\fi
\noindent\textit{Source reference.} Original lemma 02NF: focused-seeds/AGP/sources/stacks-snapshot-20261011/tags.txt, source lines 3478--3478. The lemma's own canonical tag, quoted so that readers can resolve the label inside the delivered text; it points at the counted statement itself and carries no proof content.\par\medskip
\ifdefined\hypertarget\hypertarget{sourcepointer-2}{}\fi
\noindent\textit{Source reference.} Original lemma 02NA: focused-seeds/AGP/sources/stacks-snapshot-20261011/tags.txt, source lines 3473--3473. Navigation pointer to a statement the counted proof refers to but does not contain; nothing load-bearing is delegated to it.\par\medskip
\ifdefined\hypertarget\hypertarget{sourcepointer-3}{}\fi
\noindent\textit{Source reference.} Original lemma 01KM: focused-seeds/AGP/sources/stacks-snapshot-20261011/tags.txt, source lines 2055--2055. Navigation pointer to a statement the counted proof refers to but does not contain; nothing load-bearing is delegated to it.\par\medskip
\ifdefined\hypertarget\hypertarget{sourcepointer-4}{}\fi
\noindent\textit{Source reference.} Original lemma 01J7: focused-seeds/AGP/sources/stacks-snapshot-20261011/tags.txt, source lines 2004--2004. Navigation pointer to a statement the counted proof refers to but does not contain; nothing load-bearing is delegated to it.\par\medskip
\begin{lemma}
\label{lemma-distinct-local-rings}
Let $X$ be an integral separated scheme.
Let $Z_1$, $Z_2$ be distinct irreducible closed subsets of $X$.
Let $\eta_i$ be the generic point of $Z_i$.
If $Z_1 \not\subset Z_2$, then
$\mathcal{O}_{X, \eta_1} \not \subset \mathcal{O}_{X, \eta_2}$
as subrings of $R(X)$.
In particular, if $Z_1 = \{x\}$ consists of one closed point $x$,
there exists a function regular in a neighborhood of $x$
which is not in $\mathcal{O}_{X, \eta_{2}}$.
\end{lemma}

\begin{proof}
First observe that under the assumption of $X$ being separated,
there is a unique map of schemes
$\Spec(\mathcal{O}_{X, \eta_2}) \to X$ over $X$
such that the composition
$$
\Spec(R(X)) \longrightarrow
\Spec(\mathcal{O}_{X, \eta_2}) \longrightarrow X
$$
is the canonical map $\Spec(R(X)) \to X$.
Namely, there is the canonical map
$can : \Spec(\mathcal{O}_{X, \eta_2}) \to X$, see
Schemes, Equation (\ref{schemes-equation-canonical-morphism}).
Given a second morphism $a$ to $X$, we have that $a$ agrees with $can$
on the generic point of
$\Spec(\mathcal{O}_{X, \eta_2})$ by assumption.
Now $X$ being separated guarantees that the subset in
$\Spec(\mathcal{O}_{X, \eta_2})$ where
these two maps agree is closed, see
Schemes, Lemma \ref{schemes-lemma-where-are-they-equal}.
Hence $a = can$ on all of $\Spec(\mathcal{O}_{X, \eta_2})$.

\medskip\noindent
Assume $Z_1 \not \subset Z_2$ and assume on the contrary that
$\mathcal{O}_{X, \eta_{1}} \subset \mathcal{O}_{X, \eta_{2}}$
as subrings of $R(X)$. Then we would obtain a second morphism
$$
\Spec(\mathcal{O}_{X, \eta_{2}}) \longrightarrow
\Spec(\mathcal{O}_{X, \eta_{1}}) \longrightarrow
X.
$$
By the above this composition would have to be equal to $can$.
This implies that $\eta_2$ specializes to $\eta_1$ (see
Schemes, Lemma \ref{schemes-lemma-specialize-points}).
But this contradicts our assumption $Z_1 \not \subset Z_2$.
\end{proof}

\end{document}
